\clearpage\SourceLeaf{342}
\begin{center}
\textbf{[THE CIRCLE.]}\par\medskip
[I.]\par\medskip
\textbf{COMPLEX NUMBERS}\par
\textbf{APPLIED TO SPACE.}
\end{center}

$\alpha$, $\beta$, $\gamma$, $\delta$ and\,so\,on are cyclic quantities, also denoting the points on the circumference of the circle of radius 1.

\Needspace{10\baselineskip}1. Area of the triangle $\alpha$, $\beta$, $\gamma$
\[
=\frac{(\beta-\alpha)(\gamma-\beta)(\alpha-\gamma)}{4\alpha\beta\gamma}\cdot i
\qquad \text{equilateral: }\frac{\BareRoot{27}}{4}.
\]
\[
\begin{aligned}
&=\frac{\alpha\alpha\beta+\beta\beta\gamma+\gamma\gamma\alpha-\alpha\beta\beta-\beta\gamma\gamma-\gamma\alpha\alpha}{4\alpha\beta\gamma}\cdot i\\
&=\tfrac{1}{4}i\cdot\left\{\frac{\beta}{\alpha}+\frac{\gamma}{\beta}+\frac{\alpha}{\gamma}-\frac{\alpha}{\beta}-\frac{\beta}{\gamma}-\frac{\gamma}{\alpha}\right\}.
\end{aligned}
\]

\Needspace{9\baselineskip}2. The area of the triangle that touches the circle at the points $\alpha$, $\beta$, $\gamma$ and whose vertices are represented by
\[
\frac{2\beta\gamma}{\beta+\gamma},\qquad \frac{2\gamma\alpha}{\gamma+\alpha},\qquad \frac{2\alpha\beta}{\alpha+\beta}
\]
is given as follows;
\[
\text{Area: }\frac{(\gamma-\beta)(\alpha-\gamma)(\beta-\alpha)}{(\gamma+\beta)(\alpha+\gamma)(\beta+\alpha)}\cdot i,
\qquad\text{equilateral: }\BareRoot{27}.
\]
